\documentclass[10pt,journal]{IEEEtran}
% --- Core IEEEtran-friendly packages (keep this list minimal) ---
\usepackage{cite}
\usepackage{amsmath,amssymb}
\usepackage{graphicx}
\usepackage{xcolor}
% Use IEEEtran's legacy Times font setup in the XeTeX preview as well.
\ifdefined\XeTeXversion
  \usepackage[OT1]{fontenc}
\fi

% Blue text marks this revision. For an accepted-text view, replace the
% definition below with \DeclareRobustCommand{\rev}[1]{#1}.
\DeclareRobustCommand{\rev}[1]{\textcolor{blue}{#1}}
\usepackage[caption=false,font=footnotesize]{subfig}

% Tables
\usepackage{booktabs}
\usepackage{tabularx} % loads array internally
\usepackage{makecell}
\usepackage{multirow}

% Reduce Underfull \hbox warnings in narrow IEEE columns (gentler than \sloppy)
\emergencystretch=1em

% tabularx helper column + thin rule used in Table~\ref{tab:thz_isac_comparison}
\newcolumntype{Y}{>{\centering\arraybackslash}X}
\newcommand{\rowrule}{\specialrule{0.22pt}{0.4pt}{0.4pt}}

% Some TeX installations define these lengths via subfig; define only if missing.
\makeatletter
\@ifundefined{subfigcapskip}{\newlength{\subfigcapskip}}{}
\@ifundefined{subfigbottomskip}{\newlength{\subfigbottomskip}}{}
\makeatother

\title{Self-Interference Reuse for Shared-Aperture Monostatic Full-Duplex THz ISAC}

\author{Sungmin~Cho,~\IEEEmembership{Member,~IEEE,}
        Sang-Rok~Moon,~\IEEEmembership{Member,~IEEE,}
        Sooyeon~Kim,
        Minkyu~Sung,~\IEEEmembership{Member,~IEEE,}
        Seung-Hyun~Cho,~\IEEEmembership{Member,~IEEE,}
        and~Ho-Jin~Song,~\IEEEmembership{Fellow,~IEEE}%
\thanks{This work was supported by the Electronics and Telecommunications
Research Institute (ETRI) grant funded by the Korean government
[26ZH1100, Study on 3D communication technology for hyper-connectivity].
(Corresponding author: Sungmin Cho.)}%
\thanks{S. Cho, S.-R. Moon, S. Kim, M. Sung, and S.-H. Cho are with the Electronics
and Telecommunications Research Institute (ETRI), Daejeon 34129, South
Korea (e-mail: sungmincho@etri.re.kr).}%
\thanks{H.-J. Song is with the Department of Electrical Engineering and
Graduate School of Semiconductor Technology, Pohang University of Science
and Technology (POSTECH), Pohang 37673, South Korea
(e-mail: hojin.song@postech.ac.kr).}%
}

\begin{document}
\maketitle
\begin{abstract}
\rev{We investigate residual self-interference (SI) reuse in a 280-GHz
shared-aperture monostatic integrated sensing and communication (ISAC)
prototype driven by two free-running lasers. The prototype transmits a
downlink waveform while receiving its echo through the same aperture.
After passive isolation, residual
transmit leakage mixes with the target echo in a square-law detector
and provides a same-acquisition reference for post-detection channel
frequency response (CFR) processing. The colocated sensing receiver
estimates the CFR using the known transmit waveform across the
occupied bandwidth. For fixed SI power, the modeled
SI--echo cross-beat power scales as $R^{-4}$ and the echo self-beat
power scales as $R^{-8}$. In a controlled indoor experiment with an
aligned static horn target, a 20-GBd 32-cross-QAM DFT-spread OFDM
waveform carried 100~Gb/s gross over approximately 1~m. The
reference-aided error vector magnitude was $-15.77$~dB, and the
hard-decision bit error rate was $3.875\%$. A 15-GBd ranging record
near 1.10~m showed higher target-to-background contrast with
regularized CFR processing than with matched filtering. These
experiments evaluate residual-SI reuse as a design option for
wideband monostatic THz ISAC.}
\end{abstract}
\begin{IEEEkeywords}
Full-duplex, integrated sensing and communication,
self-interference reuse, square-law detection, terahertz.
\end{IEEEkeywords}

\section{Introduction}

Sixth-generation (6G) networks are expected to combine very-high-rate
connectivity with accurate localization and environmental perception for
applications such as vehicle-to-everything (V2X), autonomous mobility,
and smart manufacturing. Integrated sensing and communication (ISAC)
supports these functions by sharing spectrum, waveforms, hardware, and
signal processing between communication and sensing
\cite{liu_integrated_2022,gonzalez-prelcic_integrated_2024}. The
terahertz (THz) band is particularly attractive for ISAC because its
large available bandwidth can support very-high-rate communication and
fine ranging
\cite{song_terahertz_2022,chaccour_seven_2022,
elbir_Terahertz_2024,zhang_terahertz_2026}. Recent experimental THz ISAC systems have demonstrated these capabilities by focusing on waveform-level integration \cite{jia_unified_2018,lyu_multichannel_2024,shi_310_2026a}.

\rev{Time-division ISAC schedules sensing and data transmission in separate
intervals. Monostatic full-duplex ISAC (FD-ISAC) allows a node to transmit
a downlink waveform and receive its echo concurrently in the same band
\cite{zhang_inband_2022,duan_frequencydomain_2024,duan_inband_2025,
zhang_cochannel_2026,du_fullduplexbased_2026,du_fullduplex_2026}.
Here, full duplex refers to concurrent downlink transmission and echo
reception. We evaluate this operation with an aligned static target in
a controlled indoor testbed. At THz frequencies, where high-gain beams
impose stringent alignment requirements \cite{sung_practical_2026},
a shared aperture gives the transmitter and sensing receiver a common
boresight \cite{wu_integrated_2025}.}

The main challenge of FD-ISAC is self-interference (SI). In-band
transmitter leakage enters the sensing receiver together with the much
weaker target echo, limiting receiver linearity and dynamic range.
Conventional FD-ISAC therefore suppresses SI in multiple stages,
combining passive antenna-domain isolation with additional RF and digital
SI cancellation (SIC)
\cite{duan_inband_2025,
du_fullduplexbased_2026,du_fullduplex_2026}.
A photonics-assisted in-band FD-ISAC transceiver employed photonic SIC and
achieved 23.6~dB of SI suppression after radar dechirping
\cite{yu_photonicsassisted_2024a}. A recent shared-aperture FD-ISAC
prototype used 30~dB of passive circulator isolation together with RF
and digital SIC, achieving 126~dB of total SI suppression
\cite{du_fullduplex_2026}. \rev{These implementations serve different receive functions and use
different bandwidths and measurement reference planes.}

\rev{We investigate the use of leakage remaining after passive isolation
as a self-homodyne reference for square-law sensing. The contribution
is the way the shared aperture supplies this reference in the 280-GHz
photonic prototype and the analysis of the conditions under which
that reference is useful. The leakage and echo mix at the detector
input. The subsequent digital processing normalizes the measured
CFR and removes its constant baseline.}

The main contributions are summarized as follows.
\begin{itemize}

\item \rev{We derive the square-law sensing model for residual-SI referencing
and examine the useful SI power and joint ISAC range under stated
receiver and propagation assumptions. The model distinguishes the
SI--echo cross-beat from the echo self-beat and describes the
same-acquisition CFR reference.}

\item \rev{We test the architecture with a 280-GHz photonic prototype in a
controlled indoor setup. The experiment evaluates downlink
communication and static monostatic ranging with the shared
waveform. The 20-GHz occupied bandwidth corresponds to a nominal
monostatic range cell of 7.5~mm.}

\end{itemize}

The remainder of this paper is organized as follows.
Section~\ref{sec:system_model} presents the proposed architecture and
signal model. Section~\ref{sec:experimental_setup} describes the
experimental setup and processing. Section~\ref{sec:experimental_results}
presents the experimental results and comparison.
Section~\ref{sec:conclusion} concludes the paper.



\section{Proposed SI-Reuse Shared-Aperture Monostatic FD-ISAC System}
\label{sec:system_model}

\begin{figure}[!t]
\centering
\includegraphics[width=\columnwidth]{fig/fig1.png}
\caption{Conceptual diagram of the proposed shared-aperture monostatic
full-duplex THz ISAC architecture. LD: laser diode; IM: intensity
modulator; UTC-PD: uni-traveling-carrier photodiode;
LNA: low-noise amplifier; ADC: analog-to-digital converter; DSP: digital
signal processing.}
\label{fig:concept}
\end{figure}

Fig.~\ref{fig:concept} shows the proposed shared-aperture monostatic
FD-ISAC architecture. A duplexing network connects the transmitter and
colocated sensing receiver to a common aperture, providing a common
boresight and passive antenna-domain isolation. \rev{Passive isolation reduces the input presented to the receiver.
Part of the leakage remains available to mix with the echo at the
square-law detector. The resulting IF waveform supports
same-acquisition CFR normalization in the digital processor.}
Section~II-A defines the waveform model, Section~II-B develops the
SI-assisted sensing and CFR normalization, Section~II-C presents the
communication model, and Section~II-D derives the joint ISAC range behavior.

\subsection{Free-Running Photonic Transmitter and Waveform}
\label{sec:tx_signal}

Two free-running lasers separated by $f_{\mathrm c}$ are photomixed in a
uni-traveling-carrier photodiode (UTC-PD), with one optical carrier
modulated by the data waveform. The complex analytic THz signal is
\begin{equation}
x(t)=\big[1+m\,d(t)\big]
e^{j[2\pi(f_{\mathrm c}+\delta f)t+\Delta\theta(t)]},
\label{eq:xt}
\end{equation}
where $d(t)$ is the normalized analytic representation of the
data-bearing intermediate-frequency (IF) waveform, satisfying
$\mathbb{E}[d(t)]=0$ and $\mathbb{E}[|d(t)|^2]=1$, $m$ is the
modulation index, $\delta f$ is the beat-frequency offset, and
$\Delta\theta(t)$ is the relative phase noise. The carrier-to-signal power ratio (CSPR) is
$\mathrm{CSPR}=1/m^2$. With the above normalization, the average
power of the modulated field is $1+m^2$ times its carrier-component
power.

\rev{For a fixed target delay and a beat-frequency offset that is
approximately constant over an acquisition, $\delta f$ contributes
a delay-dependent phase term. Differential phase noise remains
between the two paths. DFT-spread OFDM (DFT-s-OFDM) permits
frequency-domain processing and can reduce peak-to-average power
ratio (PAPR) \cite{han_thz_2024}. In the monostatic architecture, the
sensing processor has access to the transmitted data across the
occupied subcarriers \cite{keskin_monostatic_2023}. The full data
waveform can therefore serve as the CFR reference, with dedicated
sensing pilots optional. Section~\ref{sec:offline_processing}
describes the reference and block processing for the recorded data.}

\subsection{SI-Assisted Sensing Branch}
\label{sec:sens_branch}

At the sensing low-noise amplifier (LNA) input, SI and the delayed target
echo are superimposed as
\begin{equation}
r(t)=\sqrt{P_{\mathrm{SI}}}\,x(t)
+\sqrt{P_{\mathrm{ec}}(R)}\,x(t-\tau)+n(t),
\label{eq:rt}
\end{equation}
where $\tau=2R/c$ and $n(t)$ is the input-referred receiver noise with
power $N=kT_0F_{\mathrm{LNA}}B$, where $k$ is Boltzmann's constant,
$T_0$ is the reference temperature, $F_{\mathrm{LNA}}$ is the LNA noise
factor, and $B$ is the receiver bandwidth. The residual SI power after
passive antenna-domain isolation is
$P_{\mathrm{SI}}=P_{\mathrm t}10^{-\mathrm{ISO}/10}$, where
$\mathrm{ISO}$ denotes the passive antenna-domain isolation.
\rev{The LNA and detector input limits constrain the allowable leakage.
The mean SI input can be compared with the LNA compression
specification. The usable operating region also depends on the
waveform peaks and on the ZBD response at its own input.}
$P_{\mathrm{SI}}$ is independent of target range, whereas the monostatic
echo power follows \cite{skolnik_introduction_2001}
\begin{equation}
P_{\mathrm{ec}}(R)=
\frac{P_{\mathrm t}G_{\mathrm t}G_{\mathrm r}
\lambda^2\sigma_{\mathrm{eff}}}
{(4\pi)^3R^4L^2},
\label{eq:pec}
\end{equation}
where $G_{\mathrm t}$ and $G_{\mathrm r}$ are the antenna gains,
$\lambda$ is the wavelength, $\sigma_{\mathrm{eff}}$ is the effective
radar cross section (RCS), and $L$ is the one-way front-end loss.
Hereafter, $P_{\mathrm t}$, $P_{\mathrm{SI}}$, and $P_{\mathrm{ec}}$
denote carrier-component powers, with $P_{\mathrm{ec}}(R)$ abbreviated
as $P_{\mathrm{ec}}$ when the range dependence is implicit.

\subsubsection{Square-Law Detection}
\label{sec:sld}

The square-law detector output is proportional to $|r(t)|^2$.
Expanding \eqref{eq:rt} gives
\begin{align}
|r(t)|^2={}&
\underbrace{2\sqrt{P_{\mathrm{SI}}P_{\mathrm{ec}}}
\,\mathrm{Re}\{x(t)x^*(t-\tau)\}}_{\text{SI--echo cross-beat}}
+\underbrace{P_{\mathrm{ec}}|x(t-\tau)|^2}_{\text{echo self-beat}}
\notag\\
&+\underbrace{P_{\mathrm{SI}}|x(t)|^2}_{\text{SI self-beat}}
+\mathcal N(t),
\label{eq:sq_expand}
\end{align}
where
\begin{equation}
\mathcal N(t)
=
2\,\mathrm{Re}\!\left\{
\left[
\sqrt{P_{\mathrm{SI}}}x(t)
+\sqrt{P_{\mathrm{ec}}}x(t-\tau)
\right]n^*(t)
\right\}
+|n(t)|^2.
\label{eq:noise}
\end{equation}

For the waveform in \eqref{eq:xt},
\begin{equation}
|x(t)|^2
=
1+m[d(t)+d^*(t)]+m^2|d(t)|^2.
\label{eq:selfbeat}
\end{equation}
For the self-beat terms, square-law detection removes the common carrier
phase and phase noise. IF bandpass
filtering then extracts the desired beat, which is linear in $m$, while
the second-order term contributes signal--signal beat interference
(SSBI), smaller by a factor of order $m$. A CSPR of 13~dB
($m=0.224$) is therefore used, making $m$ small enough to suppress SSBI
while keeping sufficient sideband power.

The SI--echo cross-beat retains the differential phase
$\phi_\tau(t)=2\pi(f_{\mathrm c}+\delta f)\tau+
\Delta\theta(t)-\Delta\theta(t-\tau)$. The frequency-related term
$2\pi(f_{\mathrm c}+\delta f)\tau$ is constant within each acquisition,
whereas the random variation is determined by the phase-noise difference
over the round-trip delay. For two independent 15-kHz-linewidth lasers,
the resulting 30-kHz beat linewidth gives a root-mean-square (rms)
differential phase variation of $2.9^\circ$ at $R=2$~m.

Accordingly, the desired analytic detector output after IF filtering is
\begin{align}
y(t)={}&
2m\left[
P_{\mathrm{SI}}
+\sqrt{P_{\mathrm{SI}}P_{\mathrm{ec}}}\,e^{j\phi_\tau}
\right]d(t)
\notag\\
&+
2m\left[
P_{\mathrm{ec}}
+\sqrt{P_{\mathrm{SI}}P_{\mathrm{ec}}}\,e^{-j\phi_\tau}
\right]d(t-\tau)
+n_{\mathrm d}(t),
\label{eq:yt}
\end{align}
\rev{Here $C_0$ and $C_\tau$ are the undelayed and delayed coefficients,
respectively. With $a=\sqrt{P_{\mathrm{SI}}}$ and
$b=\sqrt{P_{\mathrm{ec}}}e^{-j\phi_\tau}$, they are
$C_0=(a+b)^*a$ and $C_\tau=(a+b)^*b$. For the unit-power waveform,
the delayed real-IF term has power $2m^2|C_\tau|^2$. Under a phase
average with $\mathbb{E}[\cos\phi_\tau]=0$, this gives
$2m^2(P_{\mathrm{SI}}P_{\mathrm{ec}}+P_{\mathrm{ec}}^2)$.
The model approximates the detector-output noise as
$N_{\mathrm d}=N_{\mathrm{d},0}+2N(P_{\mathrm{SI}}+P_{\mathrm{ec}})$
in input-referred power-product units. Applying the carrier--noise
term to a specific receiver requires accounting for the accepted
RF image bands and noise-equivalent bandwidth. The fixed term
represents the detector and receiver floor in the same units.}

\subsubsection{SI-Normalized CFR and Range Estimation}
\label{sec:range_est}

\rev{When inter-block contributions are negligible, the block Fourier
transform of \eqref{eq:yt} gives the following occupied-subcarrier model.}
\begin{equation}
Y_k
=
2mA_k\left(C_0+C_\tau e^{-j2\pi f_k\tau}\right)D_k
+\mathcal V_k,
\label{eq:Yk}
\end{equation}
where $D_k$ and $Y_k$ are the transmitted and received spectra,
$A_k$ is the combined detector, IF, and digitizer response, and
$\mathcal V_k$ is noise. Since $D_k$ is known at the colocated sensing
processor, the CFR is obtained by zero-forcing division \cite{sturm_waveform_2011}:
\begin{equation}
\widehat H_k
=
\frac{Y_k}{D_k}
\simeq
2mA_k\left(C_0+C_\tau e^{-j2\pi f_k\tau}\right).
\label{eq:cfr}
\end{equation}

\rev{An ideal frequency-resolved reference
$\widehat H_{\mathrm{SI},k}=2mA_kC_0$ yields the normalized CFR}
\begin{equation}
\widetilde H_k
=
\frac{\widehat H_k}{\widehat H_{\mathrm{SI},k}}-1
=
\frac{C_\tau}{C_0}e^{-j2\pi f_k\tau}.
\label{eq:sinorm}
\end{equation}
\rev{In this ideal expression the common response $A_k$ cancels. For
coefficients that are effectively constant during the acquisition,
the phase of $C_\tau/C_0$ contributes a frequency-independent offset.
The experiment uses a scalar reference as described in
Section~\ref{sec:offline_processing}. After phase unwrapping, let
$\mu$ denote the fitted phase slope with respect to frequency.}
\begin{equation}
\widehat\tau=-\frac{\mu}{2\pi},
\qquad
\widehat R=\frac{c\widehat\tau}{2}.
\label{eq:slope}
\end{equation}
A constant hardware delay offset is removed by a one-time range
calibration. For multiple targets, an inverse Fourier transform of
\eqref{eq:sinorm} provides the delay profile
\cite{sturm_waveform_2011}. A target-off reference can further suppress
static clutter, while very short-delay echoes overlapping the SI
response require additional separation.


\subsubsection{Sensing SINR}

Using the phase-averaged target power and coherent CFR processing gain
$G_{\mathrm p}$, the sensing SINR is
\begin{equation}
\gamma_{\mathrm{sens}}(R)
=
\eta_{\mathrm d}G_{\mathrm p}
\frac{
2m^2\!\left[
P_{\mathrm{SI}}P_{\mathrm{ec}}
+
P_{\mathrm{ec}}^2
\right]
}{
N_{\mathrm{d},0}
+
2N\!\left[
P_{\mathrm{SI}}+P_{\mathrm{ec}}
\right]
+
\kappa m^4P_{\mathrm{SI}}^2
},
\label{eq:sinr_sens}
\end{equation}
where $\eta_{\mathrm d}\leq1$ is the waveform-inversion efficiency and
$\kappa$ is the effective in-band SSBI coefficient after IF filtering
and digital signal processing. Without SI, \eqref{eq:sinr_sens} reduces to
\begin{equation}
\gamma_{\mathrm{sens}}^{\mathrm{w/o\,SI}}(R)
=
\eta_{\mathrm d}G_{\mathrm p}
\frac{
2m^2P_{\mathrm{ec}}^2
}{
N_{\mathrm{d},0}+2NP_{\mathrm{ec}}
},
\label{eq:sinr_nosi}
\end{equation}
where echo-generated SSBI is negligible for a weak echo. Since
$P_{\mathrm{ec}}\propto R^{-4}$, the useful echo-only term scales as
$P_{\mathrm{ec}}^2\propto R^{-8}$.

With SI reuse, square-law detection introduces the SI--echo cross-beat
term $P_{\mathrm{SI}}P_{\mathrm{ec}}$. For a weak echo
($P_{\mathrm{SI}}\gg P_{\mathrm{ec}}$), when detector noise dominates
the denominator, \eqref{eq:sinr_sens} becomes
\begin{equation}
\gamma_{\mathrm{sens}}
\simeq
\eta_{\mathrm d}G_{\mathrm p}
\frac{
2m^2P_{\mathrm{SI}}P_{\mathrm{ec}}
}{
N_{\mathrm{d},0}
}.
\label{eq:sinr_weak}
\end{equation}
Because $P_{\mathrm{SI}}$ is independent of target range, the
SI-assisted sensing SINR scales as $R^{-4}$. Under the same
detector-noise-limited condition, the echo-only SINR scales as
$R^{-8}$. \rev{These scaling laws describe the detector-noise-limited regime. The
reference field in the prototype is supplied by residual leakage
inside the shared-aperture receiver.}

\subsection{Communication Branch}
\label{sec:comm_branch}

The one-way received communication power is
\begin{equation}
P_{\mathrm{rx}}(R)=
\frac{P_{\mathrm t}G_{\mathrm t}G_{\mathrm r}\lambda^2}
{(4\pi)^2R^2L^2}.
\label{eq:prx}
\end{equation}
Using the same square-law receiver model,
\begin{equation}
\gamma_{\mathrm{comm}}(R)=
\frac{2m^2P_{\mathrm{rx}}^2(R)}
{N_{\mathrm{d},0}+2NP_{\mathrm{rx}}(R)
+\kappa m^4P_{\mathrm{rx}}^2(R)}.
\label{eq:sinr_comm}
\end{equation}
Unlike the sensing branch, the one-way communication link contains no
SI--echo cross-beat; its desired component is the carrier--sideband
self-beat, with no coherent sensing-processing gain.


\subsection{ISAC Range Analysis}
\label{sec:isac_range}

The communication and sensing ranges are defined by their respective
SINR thresholds $\gamma_{\mathrm c}$ and $\gamma_{\mathrm{th}}$
\cite{sturm_waveform_2011}. The joint ISAC range is therefore
\begin{equation}
R_{\max}
=
\min\!\left(
R_{\max}^{\mathrm{comm}},
R_{\max}^{\mathrm{sens}}
\right).
\label{eq:rmax}
\end{equation}
Table~\ref{tab:simulation_parameters} summarizes the nominal parameters
used for the closed-form analysis.

\begin{table}[!t]
\caption{Parameters for the Closed-Form Range Analysis}
\label{tab:simulation_parameters}
\centering
\footnotesize
\renewcommand{\arraystretch}{1.07}
\setlength{\tabcolsep}{2.5pt}
\begin{tabular}{@{}
>{\raggedright\arraybackslash}p{0.46\columnwidth}
>{\centering\arraybackslash}p{0.24\columnwidth}
>{\raggedleft\arraybackslash}p{0.24\columnwidth}
@{}}
\toprule
Parameter & Symbol & Value \\
\midrule
Carrier / IF frequency
    & $f_{\mathrm c},f_{\mathrm{IF}}$
    & $280/11$~GHz \\
Waveform / modulation
    & --
    & DFT-s-OFDM / 32-cross-QAM \\
Symbol rate / bandwidth
    & $R_{\mathrm s},B$
    & $20$~GBd / $20$~GHz \\
Gross line rate
    & $R_{\mathrm b}$
    & $100$~Gb/s \\
Bandwidth-limited range cell
    & $\Delta R$
    & $7.5$~mm \\
Transmit power
    & $P_{\mathrm t}$
    & $-10$~dBm \\
CSPR / modulation index
    & $\mathrm{CSPR},m$
    & $13$~dB / $0.224$ \\
Beat linewidth
    & $\Delta\nu_{\mathrm b}$
    & $30$~kHz \\
TX/RX antenna gain
    & $G_{\mathrm t},G_{\mathrm r}$
    & $33$~dBi \\
Front-end loss per pass
    & $L$
    & $2.3$~dB \\
Passive antenna-domain isolation
    & $\mathrm{ISO}$
    & $25$~dB \\
SI power at LNA input
    & $P_{\mathrm{SI}}$
    & $-35$~dBm \\
LNA input compression point
    & $\mathrm{IP}_{\mathrm{1dB}}$
    & $-20$~dBm \\
Effective RCS\textsuperscript{a}
    & $\sigma_{\mathrm{eff}}$
    & $-10$~dBsm \\
LNA gain / noise figure
    & $G_{\mathrm{LNA}},\mathrm{NF}_{\mathrm{LNA}}$
    & $13/8$~dB \\
RF noise power
    & $N$
    & $-63.0$~dBm \\
Detector noise floor\textsuperscript{b}
    & $N_{\mathrm{d},0}$
    & $-89.0$~dB(mW$^2$)\\
Effective in-band SSBI coefficient\textsuperscript{c}
    & $\kappa$
    & $0.06$ \\
Coherent processing gain\textsuperscript{d}
    & $G_{\mathrm p}$
    & $30.1$~dB \\
Waveform-inversion efficiency\textsuperscript{d}
    & $\eta_{\mathrm d}$
    & $1$ \\
Required communication SINR\textsuperscript{e}
    & $\gamma_{\mathrm c}$
    & $15.75$~dB \\
Required sensing SINR\textsuperscript{f}
    & $\gamma_{\mathrm{th}}$
    & $13.2$~dB \\
\bottomrule
\end{tabular}
\vspace{2pt}
\begin{minipage}{\columnwidth}
\scriptsize
\textsuperscript{a}Reference value; varied in
Fig.~\ref{fig:isac_range}.\\
\textsuperscript{b}Equivalent square-law detector noise floor referred
to the LNA input.\\
\textsuperscript{c}Obtained from deterministic 20-GBd DFT-s-OFDM
simulation and rounded to 0.06.\\
\textsuperscript{d}$G_{\mathrm p}$: 1024-point coherent processing;
ideal waveform inversion.\\
\textsuperscript{e}\rev{Nominal additive white Gaussian noise (AWGN) threshold for
quasi-Gray 32-cross-QAM at a pre-forward-error-correction (FEC) bit
error rate (BER) of $2\times10^{-2}$. The experimental BER is
evaluated with the archived mapping.}\\
\textsuperscript{f}Single-bin post-processing threshold for a
nonfluctuating target in AWGN with $P_{\mathrm D}=0.9$ and
$P_{\mathrm{FA}}=10^{-6}$ \cite{skolnik_introduction_2001}.
\end{minipage}
\end{table}

\begin{figure}[!t]
\centering
\includegraphics[width=0.80\columnwidth]{fig/fig2.png}
\caption{Sensing SINR versus LNA-input SI power at $R=1$~m and
$\sigma_{\mathrm{eff}}=-10$~dBsm. The dashed line denotes the
weak-echo approximation in \eqref{eq:sinr_weak}.}
\label{fig:sinr_si}
\end{figure}

\begin{figure}[!t]
\centering
\includegraphics[width=\columnwidth]{fig/fig3.png}\\[-5pt]
\hspace*{0.1\columnwidth}\makebox[0pt]{\footnotesize (a)}%
\hspace*{0.5\columnwidth}\makebox[0pt]{\footnotesize (b)}%
\setlength{\abovecaptionskip}{2pt}
\caption{Communication, sensing, and joint ISAC ranges versus effective
RCS for (a) $P_{\mathrm t}=-10$~dBm and (b) $P_{\mathrm t}=0$~dBm.}
\label{fig:isac_range}
\end{figure}

Using these parameters, Fig.~\ref{fig:sinr_si} shows the sensing SINR
versus residual SI power at $R=1$~m and
$\sigma_{\mathrm{eff}}=-10$~dBsm. As $P_{\mathrm{SI}}$ increases beyond $P_{\mathrm{ec}}$, the
SI--echo cross-beat strengthens and the sensing SINR initially increases
approximately in proportion to $P_{\mathrm{SI}}$, as described by
\eqref{eq:sinr_weak}. At higher SI
levels, SI-induced noise and SSBI limit this gain and eventually
degrade the SINR. The nominal $P_{\mathrm{SI}}=-35$~dBm lies within
the resulting useful SI operating region, \rev{providing approximately 15~dB of modeled sensing SINR gain over
echo-only square-law detection. The high-SI end of the sweep is an
extrapolation of the detector model. With 13~dB of LNA gain and
1--2~dB of assumed interstage loss, an SI carrier near
$-32$ to $-31$~dBm at the LNA input would reach the reported
$-20$-dBm ZBD input criterion on an average carrier-power basis.}

Fig.~\ref{fig:isac_range} shows how this SINR gain translates into
joint ISAC coverage. At $P_{\mathrm t}=-10$~dBm and
$\sigma_{\mathrm{eff}}=-10$~dBsm, SI reuse increases the sensing range
to 1.6~m, nearly twice that of echo-only square-law sensing. The
SI-assisted sensing range therefore exceeds the 1.0-m communication
range, shifting the joint ISAC coverage from sensing-limited to
communication-limited operation.

This shift also changes the transmit-power scaling of the joint ISAC
range. In the weak-echo, detector-noise-limited regime, the
SI-assisted sensing range and the one-way communication range both
scale as $\sqrt{P_{\mathrm t}}$. \rev{Under the assumed detector model, a 10-dB increase in
$P_{\mathrm t}$ changes the predicted joint ISAC range from about
1.0 to 3.2~m. Realizing this operating point would require adequate
LNA and detector headroom at the higher leakage power.}

\section{Experimental Setup}
\label{sec:experimental_setup}

\begin{figure*}[!t]
\centering
\setlength{\subfigcapskip}{-10pt}
\setlength{\subfigbottomskip}{2pt}

\subfloat[]{%
  \includegraphics[width=0.87\textwidth,trim=0 1mm 0 0,clip]
  {fig/fig4a.png}%
  \label{fig:exp_config}%
}\\[-2pt]

\subfloat[]{%
  \includegraphics[width=0.9\textwidth,trim=0 4mm 0 0,clip]
  {fig/fig4b.png}%
  \label{fig:exp_photo}%
}\\[-3pt]

\subfloat[]{%
  \includegraphics[width=0.21\textwidth]{fig/fig4c.png}%
  \label{fig:exp_optical}%
}\hspace{2pt}%
\subfloat[]{%
  \includegraphics[width=0.69\textwidth]{fig/fig4d.png}%
  \label{fig:exp_dsp}%
}

\caption{Experimental implementation of the proposed SI-reuse
shared-aperture monostatic full-duplex THz ISAC system with a
280-GHz carrier: (a) system configuration; (b) testbed photograph;
(c) optical spectrum at the UTC-PD input; and (d) \rev{two-channel DSO acquisition with IF spectra of the communication
(C1) and sensing (C2) channels.} MZM: Mach--Zehnder
modulator; AWG: arbitrary waveform generator; OC: optical coupler;
WSS: wavelength-selective switch; EDFA: erbium-doped fiber amplifier;
VOA: variable optical attenuator; ZBD: zero-bias Schottky diode.}
\label{fig:exp}
\end{figure*}

Fig.~\ref{fig:exp}(a) and (b) show the 280-GHz proof-of-concept
testbed for simultaneous downlink communication and monostatic
sensing. The system consists of a free-running photonic THz
transmitter, a shared-aperture monostatic front end, and square-law
receivers based on zero-bias Schottky diodes (ZBDs). A single
data-bearing waveform is used for both functions, while the residual
transmitter leakage through the shared-aperture front end is reused as
the sensing reference.

\subsection{ISAC Waveform and Photonic THz Transmitter}

\rev{The archived waveform uses a pseudorandom binary sequence of order
15 (PRBS15), 32-cross-QAM, and DFT-s-OFDM blocks with 1024 occupied
subcarriers. The archived bit labels follow the transmitter's
binary indexing of the constellation points. The occupied data
subcarriers also provide the bandwidth for ranging.}

An arbitrary waveform generator (AWG, Keysight M8194A) generated the
waveform at symbol rates up to 20~GBd with an 11-GHz IF. At 20~GBd, the 20-GHz occupied bandwidth provides a
100-Gb/s gross line rate and a bandwidth-limited monostatic range cell
of $\Delta R=c/(2B)=7.5$~mm. After IF amplification, the waveform
drove a Mach--Zehnder modulator (MZM, iXBlue MXAN-LN-40) with a
13-dB CSPR.

Two independent free-running laser diodes
(LDs, Pure Photonics PPCL500), each with a 15-kHz Lorentzian
linewidth, were operated at 192.41 and 192.13~THz. LD~\#1 was
modulated by the MZM and combined with the continuous-wave LD~\#2.
A wavelength-selective switch (WSS) selected the two optical carriers
and the lower modulation sideband [Fig.~\ref{fig:exp}(c)], avoiding
double-sideband power fading under laser-frequency drift.

The optical signal was amplified and adjusted by an erbium-doped
fiber amplifier (EDFA) and a variable optical attenuator (VOA) before
photomixing in the UTC-PD (NTT IOD-PMJ-13001). The generated THz
carrier was centered at 280~GHz with a maximum output power of
$-10$~dBm. At 20~GBd, the data-bearing THz spectrum occupies
approximately 259--279~GHz.

\subsection{Shared-Aperture FD-ISAC Front End}
\label{sec:setup_core}

The UTC-PD output was connected to a shared-aperture monostatic front
end comprising a custom dual-circular-polarization (dual-CP) duplexer
and a common corrugated horn. The duplexer converted the transmitted
signal to right-hand circular polarization (RHCP) and routed it to the
horn with a measured insertion loss of 2~dB
\cite{shu_wideband_2019}. Finite isolation between the transmit and
sensing ports produced the residual SI used by the proposed
architecture.

Two TPX lenses (Thorlabs TPX100) collimated and refocused the 1-m THz
beam. A 25-dBi horn at the remote end served as both the communication
receive antenna and the sensing target. The reflected signal returned
to the common aperture with left-hand circular polarization (LHCP) and
was routed to the sensing port, where it combined with the residual SI
at the LNA input. The experimental link model used a 33-dBi effective
antenna gain and a $-7$-dBsm effective RCS.

\rev{The measured passive antenna-domain isolation averaged 23~dB over
259--280~GHz. For approximately $-10$~dBm of transmit power, this
gives an estimated mean residual SI of $-33$~dBm at the sensing LNA
input. With 13~dB of typical LNA gain and 1--2~dB of assumed
interstage loss, the estimated mean ZBD input is approximately
$-22$ to $-21$~dBm. The corresponding average margins to the
reported $-20$-dBm compression points are about 13~dB at the LNA
and 1--2~dB at the ZBD. These estimates concern mean power.
Waveform peaks and the detector's square-law range require
separate characterization.}

The superimposed SI and target echo were amplified by a THz LNA
(VDI WR3.4AMP-LN) and square-law detected by a ZBD
(VDI WR3.4ZBD). As described by \eqref{eq:sq_expand} and
\eqref{eq:yt}, the square-law detector generated the SI--echo cross-beat
in the IF band. After DC blocking and IF amplification, the sensing
signal was digitized on channel C2 of a digital storage oscilloscope
(DSO, Keysight UXR0404A) at 256~GSa/s.

The remote communication receiver used the same LNA--ZBD
direct-detection chain and was digitized on channel C1.
\rev{The two channels were recorded by the same DSO
[Fig.~\ref{fig:exp}(d)]. The archived two-channel records contain
the C1 and C2 samples and their common stored transmit reference.}

\begin{figure}[!t]
\centering
\includegraphics[width=1\columnwidth]{fig/fig5.png}
\caption{Measured and signal-to-noise ratio (SNR)-limited error vector
magnitude (EVM) versus
symbol rate over the 1-m link with a 280-GHz carrier, including 32-cross-QAM
constellations at 2 and 20~GBd. \rev{The horizontal line is a nominal AWGN benchmark for quasi-Gray
labeling.}}
\label{fig:evm_symbol_rate}
\end{figure}

\subsection{Offline ISAC Processing}
\label{sec:offline_processing}

\rev{The complete stored transmit waveform is the common reference.\footnote{%
\rev{The archived records analyzed here include an optional
superimposed Zadoff--Chu component with a nominal power weight
of 0.20. It is included in the stored reference.}}
The archived 15-GBd ranging record contains 15 DFT-s-OFDM blocks,
each with 1024 occupied subcarriers and an 8192-sample FFT. The
checked 20-GBd record contains ten 6144-sample FFT blocks with the
same number of occupied subcarriers. Both waveforms used no cyclic
prefix ($N_{\mathrm{CP}}=0$). The DSO records were downconverted
from the IF and resampled from 256 to 120~GSa/s before block
processing. The regularized sensing CFR reprocessing uses
rectangular FFT windows. At 1.10~m, the geometric round-trip delay
corresponds to approximately 880 samples at 120~GSa/s, or
$10.7\%$ of the 15-GBd FFT window. With alignment to the SI reference,
the delayed echo can contain samples from the preceding transmit
block. Equation~\eqref{eq:Yk} neglects this inter-block term.}

\rev{For communication, the synchronized C1 waveform is demodulated
with frequency-domain equalization. RMS EVM and the hard-decision
pre-FEC BER are evaluated against the stored transmit reference.
We use $1/\mathrm{EVM}_{\mathrm{rms}}^2$ as an EVM-derived
communication metric, with $\mathrm{EVM}_{\mathrm{rms}}$ expressed
in linear units.}

For sensing, the transmit reference and C2 recording are jointly
processed across the blocks to reduce sensitivity to block-dependent
spectral variations. The regularized CFR is estimated as
\begin{equation}
\widehat H_k=
\frac{\sum_{b=1}^{N_b}Y_{b,k}D_{b,k}^{*}}
{\sum_{b=1}^{N_b}|D_{b,k}|^2+\varepsilon\overline S_{DD}},
\qquad \varepsilon=10^{-3},
\label{eq:cfr_reg}
\end{equation}
where \rev{$N_b$ is the number of jointly processed
blocks ($N_b=15$ for the archived 15-GBd ranging record),}
$D_{b,k}$ and $Y_{b,k}$ are the transmitted and received complex
spectra of block $b$, respectively, and $\overline S_{DD}$ is the
mean accumulated transmit spectral power over the occupied subcarriers.

\rev{Equation~\eqref{eq:sinorm} describes frequency-resolved SI
normalization in the ideal model. The experiment uses the complex
mean of the occupied CFR bins as a scalar reference. Division by
this reference removes a common complex scale. Subtraction of
unity removes the constant CFR baseline after detection. This
subtraction acts on the zero-delay response of the occupied-bin
aperture, while frequency selectivity in $A_k$ can remain.
Uniform coherent integration over the occupied bins yields the
range profile. A known-range calibration accounts for the fixed
offset between the SI and echo delay references.}

\rev{The sensing profile is characterized by a target-to-background
excess statistic. In linear power units, let $P_{\mathrm{pk}}$
be the target-peak power and $P_{\mathrm{fl}}$ the median
off-target floor.}
\begin{equation}
\rev{Q_{\mathrm{prof,dB}}}
=10\log_{10}
\left(\frac{P_{\mathrm{pk}}}{P_{\mathrm{fl}}}-1\right).
\label{eq:sinr_profile}
\end{equation}
The target peak is the maximum-power bin within $\pm40$~mm of the
nominal range. The median floor is evaluated over 0.5--2.0~m,
excluding a $\pm2\Delta R$ guard interval around the detected target.
\rev{The excess statistic is defined for $P_{\mathrm{pk}}>P_{\mathrm{fl}}$.
Its value depends on the search window, background statistics,
and processing sidelobes. The analytical SINR in
\eqref{eq:sinr_sens} uses a mean modeled noise and interference power
and requires separate calibration for comparison.}

\section{Experimental Results and Discussion}
\label{sec:experimental_results}

\begin{figure}[!t]
\centering
\includegraphics[width=0.72\columnwidth]{fig/fig6.png}
\caption{Matched-filter and SI-normalized CFR range profiles at
15~GBd for a target near 1.10~m. ROI: region of interest.}
\label{fig:cfr_profile}
\end{figure}

\subsection{Communication Performance}

\rev{Fig.~\ref{fig:evm_symbol_rate} shows measured EVM and an SNR-limited
benchmark over the 1-m link. The 20-GBd constellation shows more
spreading than the 2-GBd constellation. In the archived 20-GBd
32-cross-QAM record, the gross line rate was 100~Gb/s. Reprocessing
the C1 samples yielded $-15.768$~dB EVM and 1984 errors over
51\,200 bits, corresponding to a hard-decision BER of $3.875\%$.
The $-15.75$-dB line in the figure refers to the nominal
quasi-Gray AWGN criterion. The experimental BER is reported using
the archived labeling.}

\subsection{SI-Normalized CFR Ranging}

\rev{Fig.~\ref{fig:cfr_profile} shows the matched-filter and SI-normalized
CFR profiles for a 15-GBd acquisition near 1.10~m. Both profiles
place the peak near 1.098~m, and the CFR profile shows higher
target-to-median-background contrast. This is a comparison of two
processing methods applied to an SI-present acquisition.}

% Author review: The submitted Fig. 7 is retained below for reference.
% Its sensing axis combines a median-floor profile statistic with
% detector-output model SINR. Revisit after harmonizing the metrics
% and tracing the paired source captures.
\iffalse
\begin{figure}[!t]
\centering
\includegraphics[width=0.87\columnwidth]{fig/fig7.png}
\caption{Measured communication and sensing SINRs versus range at
15~GBd, together with the detector-output model predictions from
\eqref{eq:sinr_sens} and \eqref{eq:sinr_comm}.}
\label{fig:sinr_range}
\end{figure}
\fi

\subsection{\rev{Interpretation of the Ranging Results}}

\rev{The measured profile statistic in \eqref{eq:sinr_profile}
characterizes the detected peak relative to the observed background.
The modeled SINR in \eqref{eq:sinr_sens} uses a different
denominator. A quantitative comparison and the application of a
joint-range threshold to the experiments require calibration of the
profile statistic and paired communication and sensing records
at each distance. The echo-only SINR is a model baseline.}

\begin{figure}[!t]
\centering
\subfloat[]{%
  \includegraphics[width=0.6\columnwidth]{fig/fig8a.png}%
  \label{fig:range_stage}}

\subfloat[]{%
  \includegraphics[width=0.73\columnwidth]{fig/fig8b.png}%
  \label{fig:range_profiles}}
\caption{Resolution-scale displacement measurement:
(a) target translation of approximately 7.5~mm checked with a digital
caliper and
(b) corresponding 20-GBd SI-normalized CFR range profiles showing a
7.5-mm range-peak shift.}
\label{fig:range_resolution}
\end{figure}

\begin{table*}[!t]
\caption{Representative millimeter-wave, terahertz,
and full-duplex ISAC demonstrations}
\label{tab:thz_isac_comparison}
\centering
\footnotesize
\setlength{\tabcolsep}{1.0pt}
\renewcommand{\arraystretch}{1.08}

\renewcommand{\tabularxcolumn}[1]{m{#1}}

\begin{tabularx}{\textwidth}{@{}
>{\raggedright\arraybackslash}m{0.105\textwidth}
*{7}{Y}
@{}}
\toprule
&
\makecell{JSTQE'23\\\cite{bai_photonicsassisted_2023}}
&
\makecell{JLT'24\\\cite{lyu_multichannel_2024}}
&
\makecell{JLT'26\\\cite{shi_310_2026a}}
&
\makecell{OE'24\\\cite{yu_photonicsassisted_2024a}}
&
\makecell{IoTJ'25\\\cite{duan_inband_2025}}
&
\makecell{WCM'26\\\cite{du_fullduplex_2026}}
&
\textbf{This work}
\\
\midrule

Sensing mode
& Monostatic
& Bistatic
& Bistatic
& Monostatic
& Monostatic
& Monostatic
& \textbf{Monostatic}
\\
\rowrule

TX/RX aperture
& Separate
& Separate
& Separate
& Separate
& Separate
& Shared
& \textbf{Shared}
\\
\rowrule

SI strategy
& --
& --
& --
& \makecell{Photonic SIC\\(23.6~dB)}
& \makecell{Passive + spatial\\+ digital SIC}
& \makecell{Passive + RF\\+ digital SIC\\(126~dB)}
& \makecell{\textbf{Passive isolation}\\\textbf{+ SI reuse}\\
\rev{+ CFR baseline}\\\rev{removal}}
\\
\rowrule

Sensing RX
& \makecell{Photonic dechirp\\+ PD}
& \makecell{Schottky mixer\\+ LO}
& \makecell{FMD direct\\detection}
& \makecell{PD + mixer after\\photonic SIC}
& \makecell{$16{\times}16$\\phased-array}
& \makecell{USRP B210\\SDR RX}
& \makecell{\textbf{ZBD direct}\\\textbf{detection}}
\\
\rowrule

Ranging
& \makecell{Dechirp\\+ fusion}
& \makecell{Matched filter\\+ fusion}
& \makecell{Weighted log-SNR\\+ path-loss model}
& Dechirp
& \makecell{OFDM range--\\Doppler}
& \makecell{Dechirp\\+ DFT}
& \makecell{\textbf{SI-normalized}\\\textbf{CFR}}
\\

\midrule\midrule

Carrier
& 38/52~GHz
& \makecell{250/275/\\300~GHz}
& 300~GHz
& 10~GHz
& 26.8~GHz
& 1.74~GHz
& \textbf{280~GHz}
\\
\rowrule

Waveform
& CE-LFM-OFDM
& LFM-16QAM
& PCS-OFDM
& CE-LFM-OFDM
& 16-QAM OFDM
& OFDM--chirp
& \makecell{\textbf{32-cross-QAM}\\\textbf{DFT-s-OFDM}}
\\
\rowrule

Bandwidth
& $4~\mathrm{GHz}\times2$
& $10~\mathrm{GHz}\times3$
& 33~GHz
& \makecell{2~GHz total\\(1.5~GHz sensing)}
& 491.52~MHz
& 40~MHz
& \textbf{20~GHz}
\\
\rowrule

Gross data rate
& 16~Gb/s
& 120~Gb/s
& 310.4~Gb/s\textsuperscript{a}
& 3.0~Gb/s\textsuperscript{b}
& N/R
& N/R
& \textbf{100~Gb/s}
\\
\rowrule

\rev{Reported range metric}
& 8.6~mm\textsuperscript{c}
& 2.5~mm\textsuperscript{c}
& N/R
& 100~mm
& 305.2~mm
& N/R
& \textbf{7.5~mm}\textsuperscript{d}
\\
\rowrule

Comm. distance
& 1.2~m
& \multirow{2}{*}{0.5~m}
& 0.3~m
& \multirow{2}{*}{1.3~m}
& N/R
& 40~m
& \multirow{2}{*}{\textbf{1.0~m}}
\\

\cmidrule(lr){1-2}
\cmidrule(lr){4-4}
\cmidrule(lr){6-7}

Sensing distance
& 3.2~m
&
& 0.3--1.0~m
&
& $\sim$85--145~m
& 5--20~m
&
\\

\bottomrule
\end{tabularx}

\vspace{1.5pt}
\parbox{\textwidth}{\scriptsize
\textsuperscript{a}Gross line rate obtained using
dual-polarization 33-GBd PCS-64-QAM and satisfying the NGMI
criterion at 0.3~m.
\textsuperscript{b}Calculated from the reported 64-QAM and
500-MBd CE-OFDM.
\textsuperscript{c}Obtained after frequency-channel fusion.
\textsuperscript{d}Bandwidth-limited monostatic range cell,
assessed through single-target displacement using range-profile
analysis inspired by \cite{lyu_multichannel_2024}.
``--'' denotes not applicable. N/R: not reported.
CE: constant envelope; DFT: discrete Fourier transform; FMD:
Fermi-level-managed barrier diode; LFM: linear frequency modulation;
NGMI: normalized generalized mutual information; PCS: probabilistic
constellation shaping; SDR: software-defined radio; SIC:
self-interference cancellation; USRP: universal software radio peripheral.
}
\end{table*}

\subsection{Resolution-Scale Displacement Measurement}

In the present shared-aperture FD-ISAC setup, the communication
receiver also serves as the sensing target. Displacement sensing is
therefore evaluated by translating this receiver along the range
direction and tracking the SI-normalized CFR peak. The symbol rate
was increased to 20~GBd, providing a 20-GHz occupied bandwidth and
a bandwidth-limited monostatic range cell of 7.5~mm.

\rev{The reported caliper reading indicates a displacement of
approximately 7.5~mm [Fig.~\ref{fig:range_resolution}(a)]. The
two plotted CFR peaks are 1012.5 and 1005.0~mm
[Fig.~\ref{fig:range_resolution}(b)]. This example evaluates
single-target displacement on the scale of one nominal range
cell. Absolute range accuracy and two-target separation
require separate experiments.}

\subsection{Comparison and Future Research Directions}
\label{sec:discussion}

Table~\ref{tab:thz_isac_comparison} compares representative
millimeter-wave/THz and FD-ISAC demonstrations. Previous
waveform-integrated ISAC systems have achieved communication rates up
to 310.4~Gb/s and millimeter-scale ranging
\cite{lyu_multichannel_2024,shi_310_2026a}, but do not explicitly
address SI-limited same-band monostatic FD operation. Recent FD-ISAC
prototypes instead suppress SI using photonic SIC
\cite{yu_photonicsassisted_2024a}, passive and spatial isolation
followed by digital SIC \cite{duan_inband_2025}, or cascaded passive,
RF, and digital SIC \cite{du_fullduplex_2026}. In particular, the
shared-aperture prototype in \cite{du_fullduplex_2026} combines
30~dB of passive circulator isolation with additional RF and digital
SIC to achieve 126~dB of total cancellation. \rev{The prototypes in the comparison use different bandwidths,
receiver functions, and hardware configurations.}

\rev{The 280-GHz prototype uses passive antenna-domain isolation,
residual leakage as the square-law reference, and post-detection
CFR normalization and baseline removal. A separate homodyne
reference with the same field at the detector would produce the
same delayed linear beat term. The distinction here is the way
the shared-aperture receiver obtains that field.}

\rev{The controlled indoor experiment uses a static, aligned horn at
the remote end for one-dimensional ranging near 1~m. This
configuration tests the proposed transmit--echo operation and
provides an initial evaluation of the receiver architecture.
Performance with independent targets, changing aspects, clutter,
and motion remains open. Increased transmit power, antenna gain,
or detector sensitivity could extend coverage, subject to the
residual SI level and the available receiver headroom.}

Future research directions include dynamic and multidimensional
sensing. The SI--echo cross-beat preserves target-dependent Doppler information of a
moving target, making range--Doppler sensing a natural extension of
the present static experiment. Such processing will require
compensation for inter-acquisition frequency and phase drift from the
free-running transmitter. Multi-antenna or multiple-input multiple-output
(MIMO) implementations could
further extend the architecture to angular estimation and
multidimensional localization. The colocated monostatic architecture
can facilitate transmit--receive timing synchronization, while direct
detection avoids separate coherent THz local-oscillator distribution
and locking. Multichannel operation would nevertheless require
channel-to-channel calibration and management of SI coupling across
the array. \rev{A future extension to simultaneous uplink, downlink, and sensing
would require separation of the incoming data signal, residual SI,
and target echoes within the receiver dynamic range.}

\section{Conclusion}
\label{sec:conclusion}

\rev{This work evaluated residual-SI reuse in a shared-aperture 280-GHz
photonic ISAC prototype. Passive isolation reduces transmit
leakage, and the residual field provides a reference for
square-law detection of the echo. At fixed SI power, the
modeled cross-beat and echo self-beat powers scale as $R^{-4}$
and $R^{-8}$, respectively. After detection, the receiver uses
the acquired CFR for normalization and baseline removal.}

\rev{In the controlled indoor setup, the 20-GBd record carried
100~Gb/s gross over approximately 1~m and yielded a hard-decision
BER of $3.875\%$ with the archived mapping. The static ranging
results show a gain in profile contrast with CFR processing and
a displacement near one nominal range cell in the reported
example. Together, these results support residual-SI referencing
as a design option for a shared-aperture photonic transmit--echo
receiver.}

\bibliographystyle{IEEEtran}
\bibliography{ref_fin}
\end{document}
